\documentclass{article}
\usepackage{mathnotes}

\title{Neural Networks}

\begin{document}

\section{Simple Example}

\subsection{Digit Classification}
Let's talk about an example. We have $28 \times 28$ pixel greyscale input image and we want to predict which digit of 0-9 it is.

Each of our pixels are represented by a value from 0 to 255 representing the intensity of the greyscale pixel, with 0 being black and 255 being white.

At the most basic level, we want to make a function

\[ F_{\beta} : M(\mathbb{Z}_{255})_{24,24} \to \mathbb{Z}_{10} \]

i.e. $F_{\beta}$ takes in a $24 \times 24$ matrix of \@{integers} between $0$ and $255$ and outputs an integer between $0$ and $9,$ the predicted digit represented by the image.

\subsection{As A Neural Network}

Let's start by defining a neuron.

\begin{definition}[Neuron]
In a neural network, a \textbf{neuron} $N$ is a computation unit that takes $n$ input values, applies an \@{affine-transformation}, then a typically non-linear activation function, and returns a single value.

That is, it is a function of the form

\[ N: \reals^n \to \reals, \quad N(\vec{x}) = \phi(\vec{w} \cdot \vec{x} + b), \]

where

\begin{itemize}
\item $n$ is the number of inputs to the neuron
\end{itemize}

\begin{itemize}
\item $\vec{w} \in \reals^n$ is a vector of weights,
\end{itemize}

\begin{itemize}
\item $b \in \reals$ is a bias value, and
\end{itemize}

\begin{itemize}
\item $\phi : \reals \to \reals$ is an activation function.
\end{itemize}
\end{definition}

\begin{definition}[Layer]
A \textbf{layer} in a neural network is a collections of \@{neurons} that operate in parallel on the same input \@{vector}.

A layer with $n_{\ell - 1}$ inputs and $n_{\ell}$ \@{neurons} defines a \@{function}

\[ F_{\ell} : \reals^{\ell - 1} \to \reals^{\ell} \]

of the form

\[ F_{\ell}(\vec{x}) = \phi_{\ell} \left ( W_{\ell} \vec{x} + \vec{b}_{\ell} \right ), \]

where:

\begin{itemize}
\item $W_{\ell} \in \reals^{n_{\ell} \times n_{\ell - 1}}$ is the weight matrix (with the $i$th row representing the weights for the inputs to the $i$th \@{neuron} in the layer),
\end{itemize}

\begin{itemize}
\item $\vec{b}_{\ell} \in \reals^{n_{\ell}}$ is the bias vector (with the $i$th entry representing the bias on the $i$th \@{neuron} in the layer,
\end{itemize}

\begin{itemize}
\item $\phi_{\ell} : \reals \to \reals $ is an activation function, applied \@{componentwise}, 
\end{itemize}

\begin{itemize}
\item each coordinate of $F_{\ell}(\vec{x})$ is the output of a single \@{neuron} in that layer.
\end{itemize}

A layer takes an input \@{vector}, applies the same \@{affine-transformation} to all \@{neurons} (via a shared weight \@{matrix} and bias \@{vector}), and then applies an \@{activation-function} to each neuron's output.

Its output is the \@{vector} of all \@{neuron} outputs in that layer, and in this way, we can view a layer as a \@{vector} of \@{neurons}.
\end{definition}

Finally we'll define neural network itself!

\begin{definition}[Neural Network]
A \textbf{neural network} is a function obtained by \@{composing} \@{finitely-many} \@{neurons} arranged in layers:

\[ F(\vec{x}) = (\phi_L \circ A_L) \circ \cdots \circ (\phi_1 \circ A_1)(\vec{x}), \]

where each

\begin{itemize}
\item $A_{\ell}(\vec{x}) = W_{\ell} \vec{x} + \vec{b}_{\ell}$ is an \@{affine-transformation},
\end{itemize}

\begin{itemize}
\item $\phi_{\ell}$ is an \@{activation-function} applied \@{componentwise},
\end{itemize}

\begin{itemize}
\item and $W_{\ell}, \vec{b}_{\ell}$ are learnable parameters.
\end{itemize}
\end{definition}

Our pixels are represented by a value from 0 to 1 that represents intensity, with 0 being

\end{document}
